% Part: first-order-logic
% Chapter: model-theory
% Section: reducts-and-expansions

\documentclass[../../../include/open-logic-section]{subfiles}

\begin{document}

\olfileid[pt]{mod}{bas}{red}
\section{Redutos e Expansões}

Frequentemente, é útil ou necessário comparar linguagens que têm
símbolos em comum, bem como !!{structure}s para essas linguagens. O
caso mais comum ocorre quando todos os símbolos de !!a{language}~$\Lang{L}$
também fazem parte de !!a{language}~$\Lang{L'}$, isto é, $\Lang{L} \subseteq
\Lang{L'}$. Uma $\Lang{L}$-!!{structure}~$\Struct{M}$ pode sempre
ser expandida a uma $\Lang{L'}$-!!{structure}, acrescentando-se interpretações
dos símbolos adicionais e mantendo-se inalteradas as interpretações dos
símbolos comuns. Por outro lado, de uma
$\Lang{L'}$-estrutura~$\Struct{M'}$ podemos obter uma
$\Lang{L}$-estrutura simplesmente ``esquecendo'' as interpretações dos
símbolos que não ocorrem em~$\Lang{L}$.

\begin{defn}
\ollabel{defn:reduct}
Suponha que $\Lang L \subseteq \Lang L'$, que $\Struct M$ seja uma
$\Lang L$-!!{structure} e que $\Struct M'$ seja uma $\Lang L'$-!!{structure}.
$\Struct M$ é o \emph{reduto} de $\Struct M'$ a $\Lang L$, e
$\Struct M'$ é uma \emph{expansão} de $\Struct M$ a $\Lang L'$ se, e somente se,
\begin{enumerate}
\item $\Domain{M} = \Domain{M'}$
\item Para toda !!{constant}~$c \in \Lang L$, $\Assign{c}{M} =
  \Assign{c}{M'}$.
\item Para toda !!{function}~$f \in \Lang L$, $\Assign{f}{M} =
  \Assign{f}{M'}$.
\item Para todo !!{predicate}~$P \in \Lang L$, $\Assign{P}{M} =
  \Assign{P}{M'}$.
\end{enumerate}
\end{defn}

\begin{prop}
\ollabel{prop:reduct}
Se uma $\Lang{L}$-!!{structure}~$\Struct{M}$ for um reduto de uma
$\Lang{L'}$-!!{structure}~$\Struct{M'}$, então, para todas as
$\Lang{L}$-!!{sentence}s~$!A$,
\[
\Sat{M}{!A} \text{ se, e somente se, } \Sat{M'}{!A}.
\]
\end{prop}

\begin{proof}
  Exercício.
\end{proof}

\begin{prob}
Demonstre \olref[mod][bas][red]{prop:reduct}.
\end{prob}

\begin{defn}
Quando temos uma $\Lang{L}$-estrutura $\Struct{M}$, e $\Lang{L'} =
\Lang{L} \cup \{P\}$ é a expansão de $\Lang{L}$ obtida acrescentando
um único $n$-ário !!{predicate}~$P$, e $R \subseteq \Domain{M}^n$
é uma relação de $n$ lugares, escrevemos $\Expan{M}{R}$ para a
expansão~$\Struct{M'}$ de~$\Struct{M}$ tal que $\Assign{P}{M'} = R$.
\end{defn}

\end{document}
